%% ============================================================
%% europasscv2025.cls — Example / Template Document
%% ============================================================
%%
%% USAGE:
%%   \documentclass[options]{europasscv2025}
%%
%% CLASS OPTIONS:
%%
%%   LANGUAGE (affects labels like "Date of birth", "Phone", etc.)
%%     english   — English labels (default)
%%     italian   — Italian labels
%%
%%   COLOR SCHEME (affects accent, section, highlight colors)
%%     color-default     — Default color scheme
%%     color-darkblue    — Dark blue
%%     color-blue        — Blue
%%     color-darkgreen   — Dark green
%%     color-green       — Green
%%     color-darkviolet  — Dark violet
%%     color-violet      — Violet
%%
%%   LAYOUT
%%     nologo      — Remove the Europass logo from the footer
%%     nototpages  — Remove the "Page X/Y" counter from the footer
%%
%%   Any other option is passed directly to the 'article' base class
%%   (e.g. 12pt, draft, etc.)
%%
%% COLOR OVERRIDES:
%%   Individual colors can be overridden after \documentclass using
%%   the following commands (value is a 6-digit HTML hex code):
%%     \ecvaccentcolor{RRGGBB}     — Main accent color (header background, language table background)
%%     \ecvtextcolor{RRGGBB}       — Body text color
%%     \ecvbordercolor{RRGGBB}     — Photo border color
%%     \ecvtitlecolor{RRGGBB}      — Name color in header
%%     \ecvinfocolor{RRGGBB}       — Personal info text color
%%     \ecvinfolinkcolor{RRGGBB}   — Links inside the header
%%     \ecvsectioncolor{RRGGBB}    — Section title color
%%     \ecvhighlightcolor{RRGGBB}  — Entry titles and skill highlights
%%     \ecvlinkcolor{RRGGBB}       — General hyperlink color
%%     \ecvtableitemcolor{RRGGBB}  — Text color inside language table
%%     \ecvtablesepcolor{RRGGBB}   — Separator color inside language table
%%
%% ============================================================

\documentclass[english, color-default]{europasscv2025}

%% ============================================================
%% BIBLIOGRAPHY SUPPORT (optional)
%% ============================================================
%% Load the europasscv2025-bibliography package if you want to 
%% include a list of publications in your CV. This package is 
%% optional and should only be loaded if needed, as it adds 
%% biblatex as a dependency.
%%
%% LOAD THE PACKAGE:
%%   \usepackage{europasscv2025-bibliography}
%%
%% LOAD ONE OR MORE .bib FILES (in the preamble, before \begin{document}):
%%   \addbibresource{publications.bib}
%%   \addbibresource{other.bib}
%%
\usepackage{europasscv2025-bibliography}
\addbibresource{europasscv2025_example.bib}

\begin{document}

%% ============================================================
%% PROFILE PICTURE (optional)
%% ============================================================
%% \ecvphoto{filename}
%%
%% Defines the profile picture shown in the header.
%% Must be called BEFORE the ecvinfo environment.
%% Accepts any image format supported by graphicx (jpg, png, pdf).
%% The image is automatically cropped to a circle.
%% For better result use an image with an aspect ration of 1:1.
%% To omit the photo, simply remove or comment out this line.
%%
\ecvphoto{photo.png}


%% ============================================================
%% PERSONAL INFORMATION (ecvinfo environment)
%% ============================================================
%% The ecvinfo environment renders the header with personal info.
%% Commands inside can be called in any order.
%% The header height adjusts automatically to fit the content.
%%
%% Available commands:
%%
%%   \ecvname{full name}
%%   \ecvdateofbirth{date}
%%   \ecvgender{gender}
%%
%%   \ecvnationality{nationality}
%%     Can be called multiple times to list multiple nationalities.
%%
%%   \ecvaddress[label]{address}
%%     Optional label defaults to "Home". Can be called multiple
%%     times for multiple addresses.
%%     Example: \ecvaddress[Home]{12 Strawberry Hill, Dublin}
%%              \ecvaddress[Work]{200 Rue de la Loi, Brussels}
%%
%%   \ecvphone[label]{number}
%%     Optional label defaults to "Mobile". Can be called multiple
%%     times for multiple numbers.
%%     Example: \ecvphone{(+353) 555 123 555}
%%              \ecvphone[Home]{(+353) 127 6689}
%%              \ecvphone[Work]{(+353) 999 888 777}
%%
%%   \ecvemail{address}
%%     Can be called multiple times for multiple email addresses.
%%
%%   \ecvwebsite{url}
%%     Can be called multiple times for multiple websites.
%%
%%   \ecvsocial{platform}{url}
%%     Can be called multiple times for different social networks.
%%     Example: \ecvsocial{GitHub}{www.github.com/smith}
%%              \ecvsocial{LinkedIn}{www.linkedin.com/in/smith}
%%
%%   \ecvim{platform}{username}
%%     Instant messaging / chat handles.
%%     Example: \ecvim{Whatsapp}{katie.smith}
%%              \ecvim{Telegram}{ksmith}
%%
\begin{ecvinfo}
    \ecvname{Katie Smith}
    \ecvgender{Female}
    \ecvdateofbirth{1 March 1975}
    \ecvnationality{Irish}

    \ecvaddress[Home]{12 Strawberry Hill, Dublin 8 Éire/Ireland}

    \ecvphone{(+353) 555 123 555}
    \ecvphone[Home]{(+353) 127 6689}
    \ecvphone[Work]{(+353) 999 888 777}

    \ecvemail{smith@kotmail.com}
    \ecvemail{another@email.com}

    \ecvwebsite{www.myhomepage.com}
    \ecvwebsite{www.another-homepage.com}

    \ecvsocial{GitHub}{www.github.com/smith}
    \ecvsocial{LinkedIn}{www.linkedin.com/in/katie-smith}

    \ecvim{Whatsapp}{katie.smith}
    \ecvim{Telegram}{ksmith}
\end{ecvinfo}


%% ============================================================
%% SECTIONS
%% ============================================================
%% \ecvsection{title}
%%
%% Renders a section title with a full-width decorative rule.
%% Use it to separate the main content areas of the CV.
%%
\ecvsection{Education and training}

%% ============================================================
%% ENTRIES
%% ============================================================
%% \ecventry[*]{key=value, ...}
%%
%% Renders a single CV entry (job, degree, course, etc.).
%% Available keys:
%%   title        — Entry title (bold, highlighted). Required.
%%   organization — Organization, company, or institution name.
%%   location     — City and/or country.
%%   date         — Date or date range (e.g. "2020--2023").
%%   description  — Free-form text below the title line.
%%   link         — Optional URL shown below the description.
%%
%% TWO VARIANTS:
%%   \ecventry{...}  — Work experience style: date and location are
%%                     displayed smaller and in body text color.
%%   \ecventry*{...} — Education style: date and location are displayed
%%                     at the same size and color as the title.
%%
\ecventry*{
    title={PhD --- ``Young People in the Construction of the Virtual University''},
    organization={Trinity College Dublin, The University of Dublin},
    location={Ireland},
    date={1997--2001}
}
\ecventry*{
    title={Bachelor of Science in Sociology and Psychology},
    organization={Trinity College Dublin, The University of Dublin},
    location={Ireland},
    date={1993--1997},
    description={
        Sociology of risk \newline
        Sociology of scientific knowledge / information society \newline
        Anthropology \newline
        E-learning and Psychology \newline
        Research methods
    }
}

\ecvsection{Work experience}

\ecventry{
    title={Independent Consultant},
    organization={National Youth Council of Ireland},
    location={3 Montague Street, Dublin 2, D02 V327, Ireland},
    date={August 2002 -- Present},
    description={Evaluation of European Commission youth training support
        measures for youth national agencies and young people.}
}
\ecventry{
    title={Internship},
    organization={European Commission, Youth Unit},
    location={DG Education and Culture, 200 Rue de la Loi, 1049 Brussels, Belgium},
    date={March 2002 -- July 2002},
    description={
        \textbf{Sector:} European institution. \newline
        Evaluating youth training programmes and the partnership between
        the Council of Europe and the European Commission. \newline
        Organizing and running a 2-day workshop on non-formal education
        for Action 5 large-scale projects. \newline
        Contributing to the steering group on training and developing
        action plans for the next 3 years.
    }
}
\ecventry{
    title={Researcher / Independent Consultant},
    organization={Council of Europe},
    location={Budapest, Hungary},
    date={Oct 2001 -- Feb 2002},
    description={In-depth qualitative evaluation of the 2-year Advanced
        Training of Trainers in Europe, using participant observations,
        in-depth interviews and focus groups.}
}


%% ============================================================
%% LANGUAGE SKILLS (ecvlanguages environment)
%% ============================================================
%% Renders a language skills table following the CEFR framework.
%%
%% \ecvmothertongue{language}
%%   Defines the mother tongue(s). Can be called multiple times.
%%
%% \ecvlanguage{name}{listening}{reading}{production}{interaction}{writing}
%%   Adds a row to the language table. Each skill level should be
%%   a CEFR level (A1, A2, B1, B2, C1, C2) or any other string.
%%   Can be called multiple times for multiple languages.
%%
%% \ecvcefrlevels
%%   Optional. Adds a note at the bottom of the table explaining
%%   the CEFR levels (A1/A2 = Basic, B1/B2 = Independent,
%%   C1/C2 = Proficient).
%%
\ecvsection{Language skills}
\begin{ecvlanguages}
    \ecvmothertongue{English}
    \ecvlanguage{French}{C1}{C2}{B2}{C1}{C2}
    \ecvlanguage{German}{A2}{A2}{A2}{A2}{A2}
    \ecvcefrlevels
\end{ecvlanguages}


%% ============================================================
%% PERSONAL SKILLS (ecvskills environment)
%% ============================================================
%% Renders a skill group with an optional title.
%%
%% \begin{ecvskills}[optional title]
%%   \ecvskill{skill name}
%%   ...
%% \end{ecvskills}
%%
%% Skills are displayed on a single line separated by ' | '.
%% The optional title is displayed in bold before the skill list.
%% Multiple ecvskills environments can follow each other under
%% the same \ecvsection.
%%
\ecvsection{Skills}

\begin{ecvskills}[Communication skills]
    \ecvskill{Team work}
    \ecvskill{Mediating skills}
    \ecvskill{Intercultural skills}
\end{ecvskills}

\begin{ecvskills}[Computer skills]
    \ecvskill{Microsoft Office}
    \ecvskill{HTML}
\end{ecvskills}


%% ============================================================
%% DRIVING LICENCES (ecvdrivinglicences environment)
%% ============================================================
%% Renders a list of driving licence categories.
%%
%% \begin{ecvdrivinglicences}
%%   \ecvdrivinglicence{category}
%%   ...
%% \end{ecvdrivinglicences}
%%
%% Each \ecvdrivinglicence call adds one category (A, B, C, etc.).
%%
\ecvsection{Driving Licence}
\begin{ecvdrivinglicences}
    \ecvdrivinglicence{A}
    \ecvdrivinglicence{B}
\end{ecvdrivinglicences}


%% ============================================================
%% PUBLICATION SECTIONS
%% ============================================================
%% Renders a section with a filtered bibliography.
%% Numbering is continuous across all \ecvpublications blocks.
%%
%% The optional [type] filters entries by biblatex entry type.
%% If [type] is omitted, all entries are printed.
%%
%%   EXAMPLE:
%%     \ecvpublications[thesis]{Thesis}
%%     \ecvpublications[article]{Journal Articles}
%%     \ecvpublications[inproceedings]{Conference Papers}
%%
\ecvpublications[thesis]{Thesis}
\ecvpublications[article]{Journal Articles}

%% To include only specific entries instead of all, use a raw 
%% refsection block for full control over which entries are 
%% included and how they are displayed.
%% Each refsection block has its own independent numbering
%% starting from [1].
%%
%% \begin{refsection}
%%     \nocite{key1, key2}   % cite specific entries
%%     \renewcommand{\section}[2]{\ecvsection{#2}}  % style the heading
%%     \printbibliography[title=My Section]
%% \end{refsection}

\end{document}